%%%%%%%%%%  scan idx 355  |  folio 342  %%%%%%%%%%
envoyant \((1,0)\) sur \(1\), et \((0,x)\) sur \(x\cdot\lambda_{-1}(N)\). Il est immédiat de voir sur (5.5.1) et (5.3.1) que c'est un \(\lambda\)-homomorphisme, et \(A\) étant intègre dans le cas considéré, on vérifie aussitôt qu'il est injectif; comme \(A\) est un \(\lambda\)-anneau, on en déduit le résultat.

Dans le \(\lambda\)-anneau \(A_N\), on peut définir des opérations \(\gamma\) (3.2). Pour \(x\in\Lambda\), la valeur de \(\gamma^p\) calculée sur l'élément \((0,x)\) de \(A_N\) est donnée par :
\[
  \gamma^p(0,x)=\sum_{i=0}^{p-1}\binom{p-1}{i}\lambda^i(N,x),
\]
c'est donc un élément de \(A\); il est défini par la relation :
\begin{equation}\tag{5.5.4}
  \gamma^p(N,x)\cdot\lambda_{-1}(N)=\gamma^p(x\cdot\lambda_{-1}(N)),
\end{equation}
et les polynômes universels qui en résultent.


\textbf{5.6.} Soit \(A\) le \(\lambda\)-anneau engendré sur \(\mathbb Z\) par les générateurs \(N,L,x\), soumis aux relations \(\lambda^k(N)=0\) pour \(k>d\), \(\lambda^k(L)=0\) pour \(k>1\). On se propose de calculer les \(\lambda^p(N+L,x)\) en fonction des \(\lambda^q(N,x)\).

Par définition :
\[
  \lambda^p(N+L,x)\cdot\lambda_{-1}(N)\cdot\lambda_{-1}(L)
  =\lambda^p(x\cdot\lambda_{-1}(N)\cdot\lambda_{-1}(L)),
\]
soit :
\[
\begin{aligned}
1&+\lambda_{-1}(N)\cdot\lambda_{-1}(L)\cdot\sum_{k\ge1}\lambda^k(N+L,x)\cdot t^k \\
 &=\lambda_t(x\cdot\lambda_{-1}(N))\circ \lambda_t(1-L) \\
 &=\frac{\lambda_t(x\cdot\lambda_{-1}(N))}{\lambda_{Lt}(x\cdot\lambda_{-1}(N))}
   \qquad\left(\text{car }\lambda_t(1-L)=\frac{\lambda_t(1)}{\lambda_t(L)}=\frac{1+t}{1+Lt}.\text{ et on applique }(2.3.2)\right) \\
 &=\frac{1+\cdots+\lambda^k(x\cdot\lambda_{-1}(N))t^k+\cdots}{1+\cdots+\lambda^k(x\cdot\lambda_{-1}(N))L^k t^k+..} \\
 &=1+\frac{\sum_{k\ge1}\lambda^k(N,x)\cdot\lambda_{-1}(N)\cdot(1-L^k)t^k}{1+\sum_{k\ge1}\lambda^k(N,x)\cdot\lambda_{-1}(N)\cdot\cdot L^k t^k}
\end{aligned}
\]
